\section{Finding one’s predominant fault }

By Fr Konrad Loewenstein, FSSP

Dowry,FSSP Periodical

N. 41, Spring 2019\footnote{\url{https://rorate-caeli.blogspot.com/2019/07/finding-ones-predominant-fault.html}}

\subsection*{Its Nature}

Each of us has a particular temperament which encompasses our whole manner of feeling, judging, sympathizing, willing, and acting. This temperament is to be perfected in each one of us by the practice of the Christian virtues. What can impede this work of perfection, and even bring each of us to our eternal ruin, is what is known as ``the Predominant Fault''. 

Fr. Garrigou-Lagrange OP describes it as ``our domestic enemy dwelling in our interior... at times it is like a crack in a wall that seems to be solid but is not so: like a crevice, imperceptible at times but deep, in the beautiful facade of a building, which a vigorous jolt may shake to the foundations.'' Like a crack, we may notice our predominant fault, but think that it is just on the surface, and does not go deep; or we may have seen it in the past but just painted it over and now we do not see it any more. Prudence dictates that, if we see a crack in a wall, we examine it and see whether it does in fact go deeper: perhaps there is a structural problem which threatens the whole edifice. 

Some examples of the predominant fault are moral weakness, sloth, gluttony, sensuality, irascibility, and pride. Our predominant fault can inform and colour our entire temperament, and compromise our predominant virtue which is, to quote Fr Garrigou-Lagrange again, ``a happy inclination of our nature'' which should develop and increase by Grace. This predominant virtue should itself determine our temperament. 

Take the example of a person who has a temperament which is passive, patient, docile, and resigned, whose predominant fault is moral weakness, whose predominant virtue is gentleness. If his predominant fault gains the ascendancy, he will become prey to human respect, moral cowardice, unreflective conformity to evil conventions and fashions, prey to excessive indulgence, and even complete loss of energy. He is no longer gentle but simply weak, although he may, with every-one else as well, regard himself as gentle, meek, good, and kind. His gentleness has been crushed, suppressed, and destroyed by moral weakness.

Similarly, some-one with a temperament which is strong may have as his predominant virtue, fortitude in confronting injustice, but as his predominant vice, anger and irascibility. The danger for this person is that he give free rein to his irascibility, so that his fortitude degenerates into unreasonable violence in words, deeds and thoughts, which does untold harm to others, but above all to himself.

What is essential is that we first recognise our predominant fault, and then combat it. If we discover moral weakness in ourselves, we must fight it by constantly asking God to make us strong, and doing our best to face up to our duties and responsibilities, and all the unpleasant challenges which life throws up continually, even if there is no-one observing us, to reprimand us for being remiss or disengaged.

If we reflect, by contrast, that we are irascible, we must work on ourselves with the Grace of God and submit our anger to rigorous control, and thereby learn gentleness and docility, even if (as in the cases of St. Ignatius Loyola and St. Francis de Sales) this may involve a labour of many years.

But if we are in this world to perfect ourselves, is this not an important work to do? – at least as important as a conscientious accomplishment of the duties of our state in life, as our daily occupations, and those works of Charity in which we may be trying to help our neighbour. Must we not love ourselves as our neighbour? and is not the true love of self the moral perfection of the self ? If our motives in all that we do are flawed by pride for example, then all that we do will be flawed; if we are weak, then we forgo many good actions which would tend in their turn to have many good consequences; if, again, we are too forceful and irascible, we bring about, in the words of the same Fr. Garrigou-Lagrange, every kind of disorder; if we are critical of others and harbour and cultivate antipathies, then we are permanently contravening Our Lord's commandment to love our neighbour.

As time passes, the predominant fault becomes a habit, and informs and colours our whole temperament, so that it becomes natural for us to feel, judge, think, and act under its influence, and it becomes hard for us to discern the presence of this fault because it has taken us over. Or if we do discern it, it will be hard for us to admit it, especially if we are proud. And if we both discern and admit the existence of this crack in the wall, we will not want to examine it, if we are morally weak, or if we fear conversion and fear to change thereby our entire lives.

Meanwhile the devil enters in with his wiles. He knows our fault. He has been working with it all our lives. He himself has painted or plastered over the crack, or helped us do so. He augments our blindness in not seeing it, our pride in not admitting it, our fear of rebuilding the whole house, if rebuild it we must.

\subsection*{How to Discover the Fault}

We have said that our predominant fault can inform and colour our entire temperament: the way that we feel, judge, sympathise, will, and act. Were we to find it, to find it out, we could gradually change our temperament for the better: to become better people; more loving towards God and neighbour; more at peace; more filled with the light of Grace; more happy; more diffusive of God's light in this world and in the next. 

It would be a great Grace, says Fr. Garrigou-Lagrange, to meet a saint who could say to us: ``This is your predominant fault; this is your predominant virtue.'' With this virtue you are to conquer the fault, and you are to inform and colour all you do, think, and say: your feelings, your desires, your whole view of life. This virtue is, as it were, the vehicle in which you are to advance through this world with generosity and single-mindedness on the path of union with God. 

It would be a Grace indeed to find such a saint, but otherwise how are we to discover our predominant fault? At the beginning of the spiritual life, at the time, for example, of our conversion, it is relatively easy to discover it. But as time passes, we become used to it and we judge all things in its light. It gains dominion over the soul, and, descending deeper into our very being, presents itself as part of our very selves. We have got used to it: indeed we identify ourselves with it, we can't stand back from it. When it takes root in us, says our learned and wise theologian, it offers a particular repugnance to being unmasked and fought, because it wishes to reign in us and over us: it hides itself, it puts on the appearance of virtue. 

Weakness clothes itself in the poor apparel of humility; pride in that of magnanimity; anger in the apparel of justice and righteous indignation. Man, the master of self-deceit, ends up by priding himself on the very defect which is his worst enemy, as though it were a virtue. If our neighbour accuses us of this very fault, we reply with complete conviction: ``My dear friend, I may have many defects, but I assure you that this is not one of them''. Even if our spiritual director mentions it, we shake our head – excuses come promptly to our mind, for the predominant fault easily excites our passions. It commands them as a master and they obey it instantly. Its fine appearance and its power drive us in the direction of impenitence. We see a remarkable instance in the case of Judas the traitor, pessimus mercator: the most terrible merchant. Thrift leads to avarice, avarice to treachery, and treachery to impenitence.

The enemy of our soul, meanwhile, who knows this fault, makes use of it to stir up trouble in and around us: to stir up strife, commotion, uproar, and unpleasantness: storms in our soul, storms in our encounters with others. In the citadel of the interior life, the predominant fault is the weak spot undefended by the virtues. Here it reigns as an enemy within the gates: concealed, disguised, and potent. The devil knows this fault, this enemy, and knows precisely where he is located. He works with him to destroy the citadel. If we ourselves do not know him, then we cannot fight him. If we cannot fight him, we have no true interior life and will make but little progress in this world.

How then do we find him? First by prayer: ``My God, what makes me resist Thy Divine Grace? Give me the strength to submit to it. Free me from my bonds, however painful that may be.'' Second, we examine our soul with merciless realism: What is the subject of my ordinary preoccupations in the morning when I awake, and when I am alone? Where do my thoughts and desires spontaneously fly? What is the ordinary cause of my sadness and of my joy? What is the general motivation of my actions and my sins? the nature of my temptations, the cause of my resistance to Grace? – particularly when it draws me away from my prayers or distracts me in them. Thirdly, what do other people criticise in me? my spiritual director, if I have one ? my family, those I live with, those who know me the best? Fourthly, how has the Holy Spirit inspired me in moments of true fervour? What does He ask me to sacrifice for love of Him?If we adopt these measures with sincerity and constancy of spirit we will come face to face with this interior enemy which enslaves us. Our Lord says in St. John's Gospel (8.34): ``Whosoever committed sin is the servant of sin.'' 

St James and St John wished to call down fire from heaven on a city that refused to receive them. But the Lord rebuked these ``Sons of Thunder'' (Boanerges, as He called them), saying: ``Ye know not of what spirit ye are. The Son of Man came not to destroy, but to save.'' (St Luke 9:55). But already at the Last Supper we see St John content only to rest his head on the Divine Heart of the Saviour, and at the end of his life he did little else, we are told, than to repeat constantly ``My little children, love one another.'' He had lost nothing of his ardour or thirst for justice, but it had become spiritualised and elevated by an extraordinary gentleness. 

\subsection*{How to Conquer the Predominant Fault}

When, with the grace of God, we have discovered our predominant fault, we must make the firm resolve to overcome it. To do so, we need a true and stable fervour of the will, or a ``promptness of the will in the service of God'', which, according to St Thomas, is the essence of true devotion. 

Now there are three principal means to overcome the predominant fault and they are: 1) prayer; 2) examination of conscience; and 3) a sanction.

\begin{enumerate}
\item Prayer: Once God has answered my prayer to show me what the predominant fault is, I should be assiduous and fervent in beseeching His help to overcome it. If I am weak, I pray ``O God my Strength! Give me strength!'' If I am irascible: ``O God my Patience! Give me patience!'' If I am sensual; ``My God and my all!'' .....and so on. The saints have prayed in the following ways: St. Louis Bertrand: ``Lord, here burn! here cut! here dry up all that hinders me from coming to Thee, that Thou mayest spare me in eternity.'' St Nicholas of Flue: “My Lord and my God! Take everything from me that hinders me from Thee! My Lord and my God! Give everything to me that will bring me to Thee! My Lord and my God! Take me from me and give me wholly to Thee.'' 

\item Examination of conscience: It is very useful to make a particular examination of conscience in the field of my predominant fault every evening: not just a general examination which is useful for every-one as part of their night prayers: to appraise their spiritual life in general; but a concentrated look at that particular weakness which has been the cause of my undoing so many times in the past. 

St Ignatius of Loyola considers it very appropriate for beginners to write down each week the number of times they have yielded to that predominant fault which seeks to reign in them like a tyrant. Fr Garrigou-Lagrange remarks: ``It is easier to laugh at this matter fruitlessly, than to apply it fruitfully.'' If we keep track of the money we spend and receive, why should we not keep a track of what we lose and gain in the spiritual field, which are losses and gains for Eternity? 

\item The Sanction: It is also very useful to impose a sanction or a penance on ourself each time we notice that we have fallen into this fault. The penance may take the form of a particular prayer, a moment of silence, or an exterior or interior mortification. This helps us to be more circumspect for the future, and makes reparation for the fault and satisfaction for the penalty owing to it. In this way many people have cured themselves, for example of blasphemy or of cursing, by obliging themselves to give alms each time they fall. In the arduous combat against the predominant fault we must take courage. We may be tempted to pusillanimity, particularly by the devil: to think that we will never be able to eradicate it, never be able to master ourselves. But we should not make peace with our faults, otherwise we will be abandoning the interior life altogether, and our one goal in this life which is perfection. God has commanded us to be perfect, so it must be possible, that is to say with His Grace. The Council of Trent declares with St Augustine: ``God never commands the impossible, but, in giving us His precepts, He commands us to do what we can, and ask for the Grace to do what we cannot.''

The other temptation to pusillanimity comes from comparing ourselves with the canonised saints. We resignedly think that this struggle against our defects is suitable only for them, so that they might reach the highest regions of spirituality and sanctity which are not for us, but reserved for them alone. And yet as we have said, did Our Lord not command us to be perfect: to love Him with all our hearts, soul, strength, and mind, and our neighbour as ourself? This, then, is the task of every-one, even if our love may never be as remarkable as that of the canonised saints with all the extraordinary talents and gifts that they received. 

\end{enumerate}

Before conquering the predominant fault, our virtues are more like good inclinations than true and solid virtues that have taken root in us. When, with God's help, we have overcome it, the virtues become firm and strong in the nourishing rays of Charity. Charity, the love of God and of souls, comes to reign in our souls through our predominant virtue. It transforms our temperament, making us more truly ourselves: ourselves without our defects, ourselves in Charity, in God.

Peace will enter the soul, and the interior joy that it brings with it, because peace is the tranquillity of order which we have re-established in our souls by our mortification, that is, by our struggle against our own evil.

We become open to God like a flower opening to the sun: no longer referring all things to ourselves, as we did when the predominant fault reigned, but referring all things to Him: thinking always of Him, living always for Him, and leading back to Him all those with whom we come into contact. God has answered our prayer to take us from ourselves and make us wholly His, whereby we have lost nothing except for our evil, and we have gained our true self, our true being in Him. Thanks be to God!